% =============================================================================
%  The Harbor Volume — Workbook
%  Part V · The Anchor Protocol (Beam A — Identity & capability)
%  SOLUTIONS
%
%  Self-contained \input fragment; reuses \wbsol from the parent workbook with a
%  local fallback. Worked answers: derivations for Trace items, design
%  discussion for Open items. Cross-references to the paper are by section name,
%  not number, so the file survives renumbering.
% =============================================================================

\providecolor{hhcobalt}{HTML}{003FB8}
\providecolor{hhteal}{HTML}{00564C}
\providecolor{hhamber}{HTML}{6B4500}
\providecolor{hhink}{HTML}{1B1712}
\providecommand{\wbtier}[1]{{\footnotesize\scshape\bfseries\textcolor{hhink}{#1}}}

\providecommand{\wbsol}[2]{%
  \par\medskip\noindent
  \makebox[0pt][r]{\textbf{V.#1}\hspace{1.2em}}%
  \ignorespaces #2\par}

\section*{Part V \quad The Anchor Protocol \quad — Solutions}
\addcontentsline{toc}{section}{Part V — The Anchor Protocol (Solutions)}

% ───────────────────────────── Check tier ──────────────────────────────────
\wbsol{1}{%
\textbf{Rule: algorithmic pinning.} The verifier records the algorithm at
\emph{issuance}, out of band, and forces the crypto library to evaluate the
signature under that pinned algorithm. It never consults the \texttt{alg} field
carried in the token. The attack works because the \texttt{alg} field is
attacker-controlled data on the wire: rewriting \texttt{alg:HS256} to
\texttt{alg:none} (strip the signature) or to a symmetric algorithm while
supplying the public key as the HMAC secret turns a key the verifier publishes
into a forging oracle. Even \emph{validating} the header against an allow-list
still lets the attacker choose \emph{which} allowed algorithm runs; if two allowed
algorithms can be confused (asymmetric verify vs.\ symmetric MAC over the same
key material), the choice is the exploit. The only safe posture is to give the
header zero authority: header trust $=0$. This is the content of the
Algorithmic-Pinning property and is exactly what the bottom row of the
algorithm-confusion figure proves: the pinned verifier admits no rewrite trace.}

\wbsol{2}{%
From parent to child, two quantities are monotone non-increasing: (i) the
capability set, $\mathrm{cap}_\text{child}\subseteq\mathrm{cap}_\text{parent}$
(in the strict-attenuation reading, $\subsetneq$); and (ii) the TTL,
$\mathrm{ttl}_\text{child}\le\mathrm{ttl}_\text{parent}$. A third invariant is
structural: the child cannot re-grant authority it does not itself hold. The
quantity that must be \emph{fresh} at every hop is the per-hop nonce $n_k$. Each
hop signs the binding tuple $(n_k,\mathrm{id}_{k-1},\mathrm{id}_k,h)$, where the
nonce defeats replay and the bound previous/next identities defeat splicing
(reusing a hop's signature inside a different chain) and substitution (changing
the message hash $h$ mid-chain).}

\wbsol{3}{%
ProVerif mechanizes Phases~1--3: symmetric pinning (authentication + injective
agreement), asymmetric identity (Ed25519 unforgeability), and multi-hop
delegation (the injective-agreement chain query). Phase~4, revocation, is
enforced at \emph{runtime} against two invariants --- filter monotonicity over a
TTL epoch and no partial reanimation --- and a property test on the cuckoo
filter's false-positive bound. It does not need the same symbolic proof because
revocation can only \emph{narrow} the set of accepted cards; it never grants
acceptance. A mechanism that strictly shrinks the admit set cannot introduce a
new authentication or escalation trace, so the existing soundness proofs are
preserved by inspection of the verification path: the revocation check is
unconditional and precedes admission.}

\wbsol{4}{%
The two senses range over different quantities. (1) \textbf{Unbounded sessions} ---
ProVerif quantifies over an unbounded number of concurrent protocol runs and
proves the security queries hold across all of them; this is delivered. (2)
\textbf{Unbounded delegation depth} --- the number of \emph{hops} in a single
chain. The modeled delegation process fixes the hop count as a structural bound
(the proof is machine-checked at depth~3, the depth realized in deployment).
Extending it to unbounded depth via recursive replication of the delegation
process is explicitly future work. Conflating the two is the trap the paper warns
against: ``unbounded sessions'' is not ``unbounded depth.''}

% ───────────────────────────── Trace tier ──────────────────────────────────
\wbsol{5}{%
Per-entry cost: $\log_2(1/\epsilon)+3$ bits. At $\epsilon=10^{-3}$,
$\log_2(10^3)=10\log_2 10\approx 9.97$, so $\approx 9.97+3\approx 13$ bits per
card --- the paper's ``$\approx 13$ bits.'' For $N=10^5$ revocations:
\[
  10^5 \times 13\ \text{bits} = 1.3\times10^6\ \text{bits}
  = \frac{1.3\times10^6}{8}\ \text{bytes} \approx 1.6\times10^5\ \text{bytes}
  \approx 160\ \text{KB}.
\]
The assumption that makes this gossip-cheap is that only \emph{fingerprints} are
shipped, not full card IDs: the version-vector deltas carry $\sim$13-bit
fingerprints per newly-revoked card, so a full filter at $10^5$ entries is
$\approx 160$\,KB and an incremental delta is a few bytes per revocation ---
trivially in memory and trivially gossiped.}

\wbsol{6}{%
Anti-entropy gossip spreads a fact epidemically: each round, every informed node
infects (in expectation) a constant fraction of the uninformed, so the informed
count grows geometrically and full coverage takes $O(\log m)$ rounds. Multiplying
the expected round count by the period $\Delta$ gives expected wall-clock
$\Delta(1+\ln m)$. For $m=10$, $\Delta=30$\,s:
\[
  30\,(1+\ln 10) = 30\,(1+2.302) = 30\times 3.302 \approx 99\ \text{s}\approx 1.7\ \text{min}.
\]
The figure's ``$\approx 2$\,minutes'' is the \emph{worst-case} headline, not the
expected value: $\Delta(1+\ln m)$ is an expectation over the random peer
selection, and the tail of the epidemic spread (the last one or two nodes to be
reached) pushes the worst case past the mean. The gap is exactly the variance of
the coupon-collector-like final phase, where the chance of picking an
already-informed peer is highest. Security-critical revocations bypass this by
immediate push to all known peers.}

\wbsol{7}{%
\textbf{The vacuity.} In the v3 model, capabilities are opaque atoms and the only
order fact is \texttt{is\_subset(c1,c1)=true}. There is no pair of distinct
capabilities $x\ne y$ with $\texttt{is\_subset}(x,y)$ derivable, so a delegated
capability can only ever \emph{equal} its parent. ``Escalation'' --- delegating
more than you hold --- is not blocked; it is \emph{inexpressible}, because there
is no larger capability to name. The no-escalation query
$\text{Accepted}(\dots\text{over-broad}\dots)\Rightarrow\text{false}$ therefore
holds trivially: its antecedent is unsatisfiable. The proof says nothing about
the protocol.\par
\smallskip
\textbf{The fix (v5).} Three additions turn the query into a real defense.
(i) A \emph{sound partial order}: \texttt{cap\_read} $\sqsubset$ \texttt{cap\_write},
both public so the attacker can construct either, with
\texttt{is\_subset(cap\_read,cap\_write)} derivable and
\texttt{is\_subset(cap\_write,cap\_read)} \emph{not}. (ii) An \emph{insider
escalation adversary}: an agent holding a valid \texttt{cap\_read} root that signs
an over-broad \texttt{cap\_write} delegation. (iii) A \emph{non-vacuity query}:
\texttt{EscalationAttempted} must be reachable, proving the attack is expressible.
ProVerif then reports Q1 (escalating delegation accepted) unreachable while Q2
(escalation attempted) reachable --- a genuine defense, not the v3 vacuum.\par
\smallskip
\textbf{The negative control.} Deleting the verifier's \texttt{is\_subset} guard
flips Q1 to false-with-trace: ProVerif exhibits the read$\to$write escalation.
This establishes that the guard is \emph{load-bearing} --- the proof passes
\emph{because} of the check, not because the model is too weak to witness the
attack. A passing proof with no counterexample under guard-removal would be
suspicious; the control is what makes the green result trustworthy.}

\wbsol{8}{%
\textbf{The attack.} Root issues $A$ a \texttt{write} card. $A$ delegates
\texttt{read} to $B$ (a legal attenuation). $B$ delegates \texttt{write} to $C$.
The chain is $A{:}\texttt{write}\to B{:}\texttt{read}\to C{:}\texttt{write}$.\par
\smallskip
\textbf{Why final-vs-root accepts it.} A verifier that checks only
$\mathrm{is\_subset}(\mathrm{cap}_\text{final},\mathrm{cap}_\text{root})$ sees
$\texttt{write}\subseteq\texttt{write}$ --- true --- and admits $C$. It never
inspects the middle hop, so it misses that $B$ held only \texttt{read} and could
not legitimately confer \texttt{write}. The middle hop is non-monotonic, but the
endpoints look consistent.\par
\smallskip
\textbf{The per-hop predicate.} Each hop must be checked against its
\emph{immediate parent}:
\[
  \mathrm{is\_subset}(\mathrm{cap}_\text{final},\mathrm{cap}_\text{mid})
  \;\wedge\;
  \mathrm{is\_subset}(\mathrm{cap}_\text{mid},\mathrm{cap}_\text{root}).
\]
On the attack chain the first conjunct is
$\texttt{write}\subseteq\texttt{read}$ --- false --- so $C$ is rejected.
\texttt{harbor\_card\_v7\_multihop\_fixed.pv} verifies exactly this and ProVerif
proves the chain rejected; \texttt{v6} mechanizes the attack against the naive
verifier. \textbf{Runtime tie-in:} the cross-harbor attenuation recompute
$\mathrm{att}_B\circ\mathrm{att}_A$ must be applied \emph{per hop}, composing
each hop's restriction in order, never collapsing to a single final-vs-root
comparison. Function composition of monotone attenuations is itself monotone only
if each factor is checked; skipping a factor breaks the chain.}

% ───────────────────────────── Open tier ───────────────────────────────────
\wbsol{9}{%
\textbf{Approach.} Replace the fixed-arity delegation process with a replicated
sub-process $!P_\text{hop}$ that reads a chain prefix off the channel, appends one
attenuated hop, and re-emits, so a single ProVerif process can generate chains of
any length. The verifier becomes a recursive walker over the reconstructed
list.\par
\smallskip
\textbf{The obstacle.} ProVerif's unbounded-session guarantee comes free because
sessions are \emph{independent} parallel runs --- the Horn-clause abstraction
quantifies over them without an inductive argument. Chain depth is different: each
hop's validity \emph{depends} on the accumulated state of all prior hops (the
running capability lower bound, the chained identity bindings). That recursion is
an inductive invariant over an unbounded data structure, which is precisely what
Horn-clause resolution does not handle natively; it tends to lose precision or
fail to terminate on recursive list predicates. The honest options are: (a) move
the depth argument to an inductive prover (Tamarin's induction, or a Coq/F$^\star$
proof of the monotone-composition lemma) and keep ProVerif for the per-hop
property; or (b) prove a depth-independent lemma --- ``if every hop satisfies the
immediate-parent predicate then the chain is attenuated'' --- once, by induction,
outside ProVerif, and let ProVerif discharge the per-hop base case. Option (b) is
the cleaner factoring and matches how the paper already separates per-hop
verification from total-depth policy.}

\wbsol{10}{%
\textbf{Taking the gossip-convergence bound.} Today the $\Delta(1+\ln m)$ figure
is borrowed from the classical anti-entropy analysis; it is not mechanized for
\emph{this} protocol's delta-shipping and version-vector reconciliation. A
credible path to \wbtier{built}: model the mesh as a probabilistic transition
system and discharge the convergence bound in a probabilistic model checker (PRISM
or Storm), proving $\Pr[\text{not converged after } k\Delta] \le f(k,m)$ for the
actual peer-selection distribution, then a separate \emph{safety} proof (TLA$^+$
or Apalache) that the version-vector merge is monotone and eventually-consistent
under reordered/duplicated deltas. The property to discharge is two-part:
\emph{safety} (no daemon ever admits a card revoked in its own converged view ---
already implied by the local gate) and \emph{liveness/quantitative} (every
revocation reaches every honest daemon within the stated time with the stated
probability). Splitting safety from the probabilistic timing keeps each tool in
its sweet spot. \emph{Alternatively}, for the constant-time comparator: rather
than trusting \texttt{subtle::ConstantTimeEq}, re-discharge branch-freedom on the
secret bytes with a dedicated constant-time checker (e.g.\ ct-verif / a
Kani-style harness asserting no secret-dependent control flow), moving it from
``audited-library trust'' to a re-proved obligation in the same source tree as the
rest of the Kani harnesses.}

\wbsol{11}{%
\textbf{Residual leakage with daily-rotating salts.} Within one salt epoch, the
hash $H(\texttt{kid},\texttt{salt}_\text{day})$ is a stable pseudonym: an
adversary cannot invert it to a \texttt{kid} (assuming $H$ is one-way and the
\texttt{kid} space has enough entropy), but can still count \emph{how many}
revocations occurred that day and observe their \emph{timing} and \emph{gossip
topology} --- which daemon revoked first, how fast it spread, and therefore which
operator is disciplining agents and at what rate. Rotation breaks
\emph{cross-day} linkage of a single card's pseudonym, so an adversary cannot
track ``the same card stayed revoked for a week.''\par
\smallskip
\textbf{What a cross-epoch adversary still infers.} (i) \emph{Volume and cadence}:
daily revocation counts and their distribution reveal an organization's incident
rate and operational rhythm even without identities. (ii) \emph{Correlation
attacks}: if the \texttt{kid} space is small or guessable, an adversary can
pre-compute $H(\texttt{kid},\texttt{salt}_\text{day})$ for each candidate once the
daily salt leaks or is itself observable, collapsing the pseudonymity --- so the
salt must be secret and high-entropy, not merely rotating. (iii) \emph{Survival
inference}: a card revoked near a day boundary appears under two pseudonyms in
quick succession; an adversary watching the rotation seam can sometimes re-link
them by timing. The honest framing the paper takes --- ``for an organization's own
daemons this is acceptable audit transparency'' --- is right for the intended
trust domain; the salted-hash scheme is a mitigation for cross-organization
gossip, not a guarantee of unlinkability, and should be stated as such.}
